% -----------------------------------------------------------------------
% Section: Methodology
% -----------------------------------------------------------------------

\subsection{Mixed-Form PINN Formulation}

The 2D incompressible Navier--Stokes equations in stress form are cast
following~\cite{rao2020pinnflow}.  A stream function $\psi$ is introduced
so that
\begin{equation}
  u = \frac{\partial\psi}{\partial y}, \qquad
  v = -\frac{\partial\psi}{\partial x},
  \label{eq:streamfunction}
\end{equation}
which satisfies the continuity equation $\partial_x u + \partial_y v = 0$
identically and eliminates one residual term.

Auxiliary stress variables $\sigma_{11}$, $\sigma_{22}$, $\sigma_{12}$
are introduced to reduce the PDE order from second to first, giving the
constitutive relations
\begin{equation}
  \sigma_{11} = 2\mu\,\partial_x u, \quad
  \sigma_{22} = 2\mu\,\partial_y v, \quad
  \sigma_{12} = \mu(\partial_y u + \partial_x v).
  \label{eq:constitutive}
\end{equation}
The momentum residuals are
\begin{align}
  f_u &= \rho\bigl(u\,\partial_x u + v\,\partial_y u + \partial_t u\bigr)
         + \partial_x p - \partial_x\sigma_{11} - \partial_y\sigma_{12},
  \label{eq:fu}\\
  f_v &= \rho\bigl(u\,\partial_x v + v\,\partial_y v + \partial_t v\bigr)
         + \partial_y p - \partial_x\sigma_{12} - \partial_y\sigma_{22}.
  \label{eq:fv}
\end{align}
The constitutive residuals enforce Eqs.~\eqref{eq:constitutive} as soft
constraints in the loss.

\subsection{Network Architecture}

The physics-informed neural network (PINN) is a fully connected
multilayer perceptron (MLP) $\mathcal{N}_\theta : \mathbb{R}^3 \to
\mathbb{R}^5$ that maps the space--time coordinate $(x, y, t)$ to the
five output fields $(\psi,\, p,\, \sigma_{11},\, \sigma_{22},\,
\sigma_{12})$.  Velocity components are recovered via automatic
differentiation of $\psi$ using Eq.~\eqref{eq:streamfunction}.

The network comprises 7 hidden layers each with 50 units and hyperbolic
tangent ($\tanh$) activations, for a total of 15\,755 trainable parameters.  Weights are initialised
with truncated Xavier (Glorot) uniform initialisation using standard
deviation $\sqrt{2/(n_\text{in}+n_\text{out})}$, truncated at $2\sigma$,
matching the TensorFlow default used in the original
implementation~\cite{rao2020pinnflow}.

\subsection{Boundary and Initial Conditions}

The spatial domain is $x\in[0,\,1.1]$~m, $y\in[0,\,0.41]$~m with a
cylinder of radius $r=0.05$~m centred at $(0.2,\,0.2)$~m.  The fluid
parameters are $\rho=1.0$~kg/m$^3$, $\mu=0.005$~Pa$\cdot$s, giving
$Re=10$.

The inlet profile is a pulsatile parabolic flow,
\begin{equation}
  u(0,y,t) = \frac{4\,U_{\max}\,y\,(H-y)}{H^2}
              \Bigl(\sin\!\Bigl(\tfrac{2\pi t}{T}+\tfrac{3\pi}{2}\Bigr)+1\Bigr),
  \label{eq:inlet}
\end{equation}
with $U_{\max}=0.5$~m/s, $H=0.41$~m, and $T=1.0$~s.  No-slip
conditions $u=v=0$ are imposed on the top, bottom, and cylinder walls.
A zero-stress (natural) outlet condition is applied at $x=1.1$~m.
Initial conditions set $u=v=p=0$ at $t=0$.

Each boundary condition and the initial condition are enforced as mean
squared error terms in the composite loss,
\begin{equation}
  \mathcal{L} = \mathcal{L}_f
    + \beta_\text{wall}\,\mathcal{L}_\text{wall}
    + \beta_\text{inlet}\,\mathcal{L}_\text{inlet}
    + \mathcal{L}_\text{outlet}
    + \mathcal{L}_\text{ic},
  \label{eq:loss}
\end{equation}
where $\beta_\text{wall}=\beta_\text{inlet}=5$ and all other weights
are unity.

\subsection{Collocation Point Sampling}

Collocation points are drawn using Latin hypercube sampling
(LHS)~\cite{rao2020pinnflow}.  A base set of 80\,000 points is drawn
uniformly over the space--time domain, supplemented by 21\,000
near-wall and wake-refinement points concentrated within two cylinder
radii and in the downstream wake region.  The total collocation set is
approximately 101\,000 points.

\subsection{Two-Phase Training}

Training proceeds in two phases.  In the first phase, the Adam
optimiser runs for 5\,000 iterations with a fixed learning rate of
$5\times10^{-4}$.  In the second phase, the L-BFGS-B algorithm (via
SciPy) is applied with up to 100\,000 function evaluations, using
history size $m=50$, maximum line-search steps of 50,
$\text{ftol}\approx\varepsilon_\text{mach}$, and $\text{gtol}=0$.
This two-phase schedule mirrors the TensorFlow training protocol of the
original paper~\cite{rao2020pinnflow}; the SciPy L-BFGS-B wrapper was
selected over the PyTorch built-in L-BFGS to faithfully reproduce the
convergence criterion.  Training converged early at approximately
68\,000 function evaluations.
